\section{Conclusion}
We present {EcoSplat}, the \textbf{first} feed-forward 3D Gaussian Splatting framework with \textbf{explicit control over the number of output Gaussian primitives} from multi-view images. EcoSplat enables flexible rendering on a wide range of end-user devices while maintaining high rendering quality. We validate the effectiveness of EcoSplat across extensive input-view settings and under various numbers of target primitives. Our experiments demonstrate that EcoSplat significantly outperforms SOTA methods using fewer primitives, and provides the freedom to dynamically control the primitive counts, offering an optimal trade-off between efficiency and rendering fidelity.

% \section{Limitations and Future Work}
% Similar to prior feed-forward 3DGS frameworks~\cite{xu2025depthsplat, chen2024mvsplat, charatan2024pixelsplat}, {EcoSplat} assumes commonly captured multi-view images with reasonable exposure variations. To further improve robustness, future work could incorporate exposure field learning~\cite{niemeyer2025nexf} and jointly train it alongside primitive prediction for more resilient reconstruction under challenging lighting conditions. It is also important to note that our current model primarily focuses on the efficient prediction of Gaussian primitives. Further efficiency gains could be achieved by optimizing the reconstruction transformer through techniques such as sparse attention~\cite{wang2025sparsevggt}, causal attention with key-value caching~\cite{streamVGGT}, or memory-based reconstruction approach~\cite{cut3r}. Incorporating these strategies could reduce computational overhead and enable faster, more scalable reconstruction without compromising rendering quality.
